% Source: https://artofproblemsolving.com/community/c3h2826137_2022_abmc_team_rounds_58__actonboxborough_math_competition
% Stable identity: aops-topic-2826137-q19
Place $A=(-3,0)$ and $B=(3,0)$. The extended sine rule gives $R=\dfrac{3}{\sin\angle ACB}$, so the condition is $30^\circ\le\angle ACB\le150^\circ$. The boundary circles have radius $6$ and centers $(0,\pm3\sqrt3)$. Above $AB$, the admissible region is inside the upper-centered disk and outside the lower-centered disk; below $AB$ the roles reverse. Thus the locus, up to boundaries of zero area, is the symmetric difference of these disks. Their center distance is $6\sqrt3$. Their intersection consists of two $60^\circ$ sectors minus two triangles of area $9\sqrt3$, hence has area $12\pi-18\sqrt3$. The desired area is $72\pi-2(12\pi-18\sqrt3)=\boxed{48\pi+36\sqrt3}$.
\begin{asy}
usepackage("amsmath");
size(7cm);
pair U=(0,3*sqrt(3)),V=(0,-3*sqrt(3)),A=(-3,0),Bp=(3,0);draw(circle(U,6));draw(circle(V,6));draw(A--Bp,blue);dot(U);dot(V);label("$A$",A,W);label("$B$",Bp,E);label("$O_+$",U,E);label("$O_-$",V,E);
\end{asy}
